\section{Introduction}
Liquid democracy is a form of direct and representative democracy, based on the concept of {\em delegation}.
Each voter has the choice of voting themselves or transferring (transitively) its vote to a trusted proxy.
Recent interest in liquid democracy, from both practical and theoretical perspectives, was sparked 
by the Pirate Party in Germany and its Liquid Feedback platform~\cite{BKN14}.
Similar initiatives have subsequently been undertaken by the 
Demoex Party in Sweden, the Internet Party in Spain, and the Net Party in Argentina.

There are many potential benefits of a transitive delegation mechanism. 
Participation may improve in quantity for several reasons. The system is easy to use and understand, induces low barriers 
to participation, and is inherently egalitarian: there is no distinction between voters and representatives; every one
is both a voter and a delegator. Participation may also improve in quality due to the flexibility to choose different 
forms of participation: voters can chose to be active participants on topics they are comfortable with or delegate on topics
they are less comfortable with. Accountability may improve due to the transparent nature of the mechanism and 
because there is a demonstrable line of responsibility between a delegated proxy and its delegators.
The quality of decision making may improve via a specialization to delegated experts
and a reduction in induced costs, such as the duplication of resources.

Our objective here is not to evaluate such claimed benefits, but we refer the reader to~\cite{Beh17, BKN14, BZ16, For02, GA15} for detailed 
discussions on the motivations underlying liquid democracy.
Rather, our focus is to quantitatively measure the performance of liquid democracy in an idealized setting.
Specifically, can equilibria in these voting mechanisms provide high social welfare?
That is, we study the {\em price of stability of liquid democracy}.

\subsection{Background}
As stated, vote delegation lies at the heart of liquid democracy. Furthermore, 
vote delegation in liquid democracy has several fundamental characteristics: optionality, retractability, 
partitionability, and transitivity. So let us begin by defining these concepts and tracing their origins~\cite{Beh17,BKN14}.

The notion of {\em optional} delegation proffers voters the choice of direct participation (voting themselves/choosing to abstain)
or indirect participation (delegating their vote). This idea dates back over a century to the work of Charles Dodgson on parliamentary 
representation~\cite{Dod1884}.\footnote{Dodgson was a parson and 
a mathematician but, as the author of  ``Alice in Wonderland'', is more familiarly known by his {\em nom de plume}, Lewis Carroll.}

Miller~\cite{Mil69} proposed that delegations be {\em retractable} and {\em partitionable}.
The former allows for delegation assignments to be time-sensitive and reversible.
The latter allows a voter to select different delegates for different policy decisions.\footnote{This option is particularly useful 
where potential delegates may have assorted competencies. For example, Alice may prefer to delegate to the Hatter on
matters concerning tea-blending but to the Queen of Hearts on matters concerning horticulture.}
 
Finally, {\em transitive} delegation is due to Ford~\cite{For02}. This allows a proxy to themselves delegate its vote {\em and} 
all its delegated votes. This concept is central to liquid democracy. Indeed, if an agent is better served by delegating her vote to 
a more informed proxy it would be perverse to prohibit that proxy from re-delegating that vote to an even more informed proxy.  
Moreover, such transitivity is necessary should circumstances arise causing the proxy to be unable to vote. It also reduces the 
duplication of efforts involved in voting.

As noted in the sixties by Tullock~\cite{Tol67}, the development of the computer opened up the possibility of large proxy voting systems.
Indeed, with the internet and modern security technologies, liquid democracy is inherently practical; see Lumphier~\cite{Lan95}.

There has been a flurry of interest in liquid democracy from the AI community. This is illustrated by the large range of 
recent papers on the topic; see, for example, \cite{Bri18,  BT18, CG17, CMM17, GKM18, 
GA15, KMP18, KR20, ZZ17}. Most directly related to our work is the game theoretic model of liquid democracy 
studied by Escoffier et al.~\cite{EGP19}. (A related game-theoretic model was also investigated by Bloembergen et al.~\cite{BGL19}.) Indeed, our motivation is an open question posed by Escoffier et al.~\cite{EGP19}:
are price of anarchy type results obtainable for their model of liquid democracy?
We will answer this question for a generalization of their model. 

%accuracy of 0-1 voting BGL19
% price of anarchy of simple game in BGL19
%supervoters GKM18, KMP18
%dynamics EGP19
%axiomatic GA15
% binary aggregation CG17, BT,
%GA, CMM-- voting in a metric space
%rankings


\subsection{Contributions}

In Section~\ref{sec:model}, we will see that vote delegation has a natural representation in terms of a 
directed graph called the {\em delegation graph}. If each agent $i$ has a utility of $u_{ij}\in [0,1]$ when agent $j$ 
votes as her delegate then a game, called the {\em liquid democracy game}, is induced on the delegation graph.
We study the {\em welfare ratio} in the liquid democracy game, which compares the social welfare of an equilibrium 
to the welfare of the optimal solution.

Pure strategy Nash equilibria need not exist in the liquid democracy game, so we focus on mixed 
strategy Nash equilibria.
Our main result, given in Section~\ref{sec:bicriteria} is that bicriteria approximation guarantees (for social welfare and rationality) 
exist in the game.
\begin{thm}\label{thm:upper}
For all $\epsilon\in [0,1]$, and for any instance of the liquid democracy game, there exists an $\epsilon$-Nash equilibrium with social 
welfare at least $\epsilon\cdot \Opt$.
\end{thm}
Theorem~\ref{thm:upper} is tight: the stated bicriteria guarantees cannot be improved.
\begin{thm}\label{thm:lower}
For all $\epsilon\in[0,1]$, there exist instances such that any $\epsilon$-Nash equilibrium has welfare at most $\epsilon\cdot \Opt +\gamma$ for any $\gamma>0$.
\end{thm}

Theorems~\ref{thm:upper} and~\ref{thm:lower} imply
that the {\em price of stability} for $\epsilon$-Nash equilibria is~$\epsilon$.
An important consequence of Theorem~\ref{thm:upper} is that strong approximation guarantees can {\em simultaneously}
be obtained for both social welfare and rationality.
Specifically, setting~$\epsilon=\frac12$ gives factor $2$ approximation guarantees for each criteria.
\begin{cor} \label{cor:upper} 
For any instance of the liquid democracy game, there exists a $\frac12$-Nash equilibrium with welfare at least~$\frac12\cdot \Opt$.
\end{cor}


